\section{Related work}
\label{sec:related}

\paragraph{Statistical practice in deep RL.}
\citet{henderson2018matters} showed that results in deep RL can change with the random seed and
with implementation details, and \citet{colas2018seeds} turned this into power analyses for the
number of seeds. \citet{jordan2020evaluating} proposed performance measures that aggregate
across environments and account for variation across runs, \citet{agarwal2021precipice} recommended interval
estimates and robust aggregates such as the interquartile mean when only a few runs are
affordable, and \citet{patterson2024empirical} gave broad guidance on empirical design. These
works treat the run as the unit and its score as given. We look inside the score: in a
procedurally generated environment a run's score also depends on which maps it is tested on and
on the sampling randomness of the evaluation episodes, and these sources call for different
remedies (more seeds, more maps, or more evaluations). We estimate them separately and show that
they differ by metric, to the point that one common summary, the acquisition verdict, is mostly
evaluation noise.

\paragraph{Generalizability theory.}
Separating the sources of variance of a measurement, and planning a design from them, is the
subject of generalizability theory \citep{cronbach1972dependability,brennan2001generalizability},
long used for test scores and ratings, with variance components estimated as in
\citet{searle1992variance}. To our knowledge it has rarely been applied to the evaluation of RL
agents; we use its standard crossed model with training seeds, maps and evaluation replicas as
facets.

\paragraph{Evaluating memory in RL.}
\citet{hausknecht2015drqn} compared a recurrent Q-network that reads one frame at a time with a
feedforward one that reads a stack of four, the classic short-memory control we also use. POPGym
\citep{morad2023popgym} found memory-free policies competitive on several of its partially
observable tasks, \citet{ni2023transformers} separated memory length from credit assignment, and
\citet{cherepanov2024memory} argued that conclusions about memory depend on the type of memory a
task requires and on the evaluation protocol. Our controls follow the same spirit for one concrete
claim: we match parameters, ablate recurrence while keeping the weights, give a short fixed
memory, and, decisively here, give memory-free agents the time step.

\paragraph{Time limits.}
\citet{pardo2018timelimits} showed that when episodes are cut by a time limit, an agent that is
not told the remaining time faces a partially observable problem, and that adding the time to the
observation helps. \citet{deverett2019interval} showed that recurrent agents trained on interval
timing tasks develop clock-like dynamics in their state. Our result is an instance of these
effects hiding inside a memory comparison on a task that is not about timing: part of what a
recurrent agent gains from its state is a clock.

\paragraph{XLand-MiniGrid.}
XLand-MiniGrid \citep{nikulin2024xland} implements MiniGrid-style worlds
\citep{chevalierboisvert2023minigrid} with objects, goals and rules in JAX and is used mainly for
meta-reinforcement learning. We use a single fixed task from it and do not study meta-learning;
our findings concern how a memory comparison on such a task should be run and read.
